\documentclass[11pt,a4paper]{article}
\usepackage[utf8]{inputenc}
\usepackage[T1]{fontenc}
\usepackage{lmodern}
\usepackage[margin=2.5cm]{geometry}
\\usepackage{graphicx}
\usepackage{booktabs}
\usepackage{longtable}
\\usepackage{array}
\\usepackage{xcolor}
\usepackage{hyperref}
\usepackage{fancyhdr}
\usepackage{titlesec}
\usepackage{enumitem}
\usepackage{tabularx}
\usepackage{colortbl}
\usepackage{float}
\usepackage{amsmath}
\usepackage{amssymb}

\definecolor{critred}{RGB}{220,38,38}
\definecolor{highorange}{RGB}{234,88,12}
\definecolor{medyellow}{RGB}{202,138,4}
\definecolor{stratblue}{RGB}{37,99,235}
\definecolor{darkbg}{RGB}{24,24,27}
\definecolor{headerblue}{RGB}{30,58,138}
\definecolor{lightgray}{RGB}{243,244,246}

\hypersetup{
  colorlinks=true,
  linkcolor=headerblue,
  urlcolor=stratblue,
  pdftitle={ICE Intelligence Briefing — KAPPA SIGINT Platform},
  pdfauthor={Samuel Wotton — Project KAPPA},
  pdfsubject={Infrastructure Vulnerability Assessment and Breach Analysis},
}

\pagestyle{fancy}
\fancyhf{}
\fancyhead[L]{\small\textsc{KAPPA SIGINT — ICE Intelligence Briefing}}
\fancyhead[R]{\small\textsc{Classified — For ICE Internal Use}}
\fancyfoot[C]{\thepage}
\fancyfoot[L]{\small ciajw.com}
\fancyfoot[R]{\small 10.0514\textdegree N, 84.2187\textdegree W}
\renewcommand{\headrulewidth}{0.4pt}
\renewcommand{\footrulewidth}{0.2pt}

\titleformat{\section}{\Large\bfseries\color{headerblue}}{\thesection}{1em}{}
\titleformat{\subsection}{\large\bfseries}{\thesubsection}{1em}{}
\titleformat{\subsubsection}{\normalsize\bfseries\color{headerblue}}{\thesubsubsection}{1em}{}

\newcommand{\priority}[2]{%
  \colorbox{#1}{\textcolor{white}{\small\textbf{#2}}}%
}

\begin{document}

\begin{titlepage}
\centering
\vspace*{2cm}

{\Huge\bfseries Intelligence Briefing}\\[0.5cm]
{\LARGE\color{headerblue} ICE Infrastructure Vulnerability Assessment}\\[0.3cm]
{\LARGE\color{headerblue} \& Breach Analysis}\\[2cm]

\rule{\textwidth}{1pt}\\[0.5cm]
{\large Prepared by the KAPPA Multi-Domain SIGINT Correlation Platform}\\[0.3cm]
{\large Project Echo — ciajw.com}\\[0.5cm]
\rule{\textwidth}{1pt}\\[2cm]

\begin{tabular}{ll}
\textbf{Classification:} & For ICE Internal Use \\
\textbf{Date:} & \today \\
\textbf{Author:} & Samuel Wotton \\
\textbf{Platform:} & KAPPA SIGINT v1.0 \\
\textbf{Location:} & Aparthotel Suites Cristina, Sabana Norte \\
\textbf{Coordinates:} & 10.0514\textdegree N, 84.2187\textdegree W \\
\textbf{Contact:} & ciajw.com \\
\end{tabular}

\vfill
{\small This document contains independently collected signal intelligence and vulnerability analysis. All findings are derived from passive collection, open-source intelligence, and publicly documented CVEs. No ICE systems were accessed or tested.}
\end{titlepage}

\tableofcontents
\newpage

\section{Executive Summary}

On March 12, 2026, the Instituto Costarricense de Electricidad (ICE) publicly disclosed a breach resulting in the exfiltration of approximately 9GB of internal email data from a locally-hosted server. Attribution has been directed at UNC2814, a China-nexus threat actor operating under the ``Gallium'' designation, using a novel command-and-control backdoor called GRIDTIDE that leverages Google Sheets API to evade traditional network detection.

This briefing presents:
\begin{enumerate}[label=\arabic*.]
  \item Two analytical frameworks for evaluating the breach attribution
  \item Verified technical evidence across 12 CVEs, 6 rootkit families, and 4 exposed SCADA protocols
  \item Specific detection indicators for GRIDTIDE C2 infrastructure
  \item A prioritized remediation plan with 8 action items
  \item Chinese satellite constellation tracking data relevant to ICE's operational environment
  \item An offer of continuous SIGINT monitoring via the KAPPA platform
\end{enumerate}

\textbf{Key Finding:} The underlying vulnerabilities are real regardless of attribution. Whether the attacker is PRC-linked, domestic, or opportunistic, ICE's SCADA systems, network edge devices, and OT infrastructure require immediate remediation.

\section{Two Analytical Frameworks}

\subsection{Theory A: Gallium Attribution as Cover}

\begin{itemize}[leftmargin=1.5em]
  \item The breach timing coincides precisely with domestic surveillance activity against individuals who documented infrastructure vulnerabilities
  \item Attributing to a Chinese APT deflects attention from internal security failures and potential misuse of monitoring infrastructure
  \item GRIDTIDE's Google Sheets C2 technique is well-documented since 2024 (TOUGHPROGRESS, Voldemort Campaign) — a convenient off-the-shelf attribution vector
  \item The 5G ban and Huawei legal battles create a pre-built geopolitical narrative
  \item Physical surveillance indicators (fiber splitters, TR-069 resets, EVOPRO room monitoring) suggest domestic actors with physical access, not remote APT operations
  \item 9GB email exfiltration from a locally-hosted server requires either insider access or persistent local network presence
\end{itemize}

\subsection{Theory B: ICE Genuinely Under Threat}

\begin{itemize}[leftmargin=1.5em]
  \item Costa Rica's critical infrastructure has real, documented vulnerabilities — SETECOM default credentials, unpatched MikroTik routers, exposed SCADA protocols
  \item UNC2814/UNC3886/UNC6384 are real, well-documented threat actors with confirmed operations across 42+ countries
  \item The Conti ransomware attack (2022) proved Costa Rica is a viable target — the ``state of war'' declaration was not theatre
  \item ICE's IT infrastructure predates modern cybersecurity standards
  \item The US \$25M cybersecurity grant acknowledges the problem is real and beyond ICE's current capacity
  \item Chinese satellite constellation provides persistent overhead coverage of Costa Rican territory
  \item Whether or not this specific breach is PRC, the vulnerabilities are real and exploitable by any motivated actor
\end{itemize}

\subsection{Convergence}

Both theories lead to identical remediation steps. The attribution question is for diplomats; the vulnerability question is for engineers. This briefing focuses on the engineering.

\section{Attack Timeline (2022--2026)}

The timeline follows a documented espionage progression: noisy ransomware (Phase 0) $\rightarrow$ infrastructure compromise (Phase 1) $\rightarrow$ silent persistent access (Phase 2).

\begin{equation}
\Psi(t) = A(t) \cdot N(t) \equiv 1
\end{equation}

As espionage amplitude $A(t)$ increases, noise $N(t)$ decreases. Phase 2 represents near-perfect operational structure.

\begin{longtable}{p{2.2cm}p{1.5cm}p{7cm}p{2cm}}
\toprule
\textbf{Date} & \textbf{Phase} & \textbf{Event} & \textbf{Severity} \\
\midrule
\endhead
2022-04-17 & Phase 0 & Conti ransomware (UNC1756) hits $\sim$30 Costa Rican institutions & \textcolor{critred}{\textbf{CRITICAL}} \\
2022-05-08 & Phase 0 & President Chaves declares national emergency — ``state of war'' & \textcolor{critred}{\textbf{CRITICAL}} \\
2022-05-31 & Phase 0 & Hive ransomware targets CCSS (EDUS + SICERE destroyed) & \textcolor{highorange}{\textbf{HIGH}} \\
2023-08-XX & Phase 1 & Budapest Convention 5G decree — Huawei/ZTE banned & \textcolor{medyellow}{\textbf{MEDIUM}} \\
2024-06-XX & Phase 1 & CISA/ZDI publish DSE855 CVEs (5947/5950/5949/5952) & \textcolor{highorange}{\textbf{HIGH}} \\
2025-06-21 & Phase 1 & Fiber splitter (NAP/Colilla) installed at Telecable distribution box & \textcolor{highorange}{\textbf{HIGH}} \\
2026-01-XX & Phase 2 & PRC bans $\sim$12 Western/Israeli cybersecurity firms (Document 79) & \textcolor{medyellow}{\textbf{MEDIUM}} \\
2026-01-25 & Phase 2 & ICE breach detected by GTIG/Mandiant — GRIDTIDE C2 identified & \textcolor{critred}{\textbf{CRITICAL}} \\
2026-01-30 & Phase 2 & Unauthorized TR-069 admin password reset on ARRIS router & \textcolor{highorange}{\textbf{HIGH}} \\
2026-03-12 & Phase 2 & ICE publicly discloses $\sim$9GB email data breach & \textcolor{critred}{\textbf{CRITICAL}} \\
2026-03-13 & Phase 2 & Media attributes ICE breach to Chinese APT ``Gallium'' & \textcolor{critred}{\textbf{CRITICAL}} \\
\bottomrule
\end{longtable}

\section{GRIDTIDE Command \& Control Architecture}

GRIDTIDE represents a novel living-off-the-cloud C2 channel that is invisible to traditional network IDS/IPS.

\subsection{Technical Profile}

\begin{tabularx}{\textwidth}{lX}
\toprule
\textbf{Attribute} & \textbf{Detail} \\
\midrule
Type & C-based backdoor \\
C2 Channel & Google Sheets API (no traditional infrastructure) \\
Encryption & AES-128-CBC, 16-byte local key file \\
Binary Masquerade & Renamed ``xapt'' (mimics apt packaging tool) \\
Persistence & systemd service at \texttt{/etc/systemd/system/xapt.service} $\rightarrow$ \texttt{/usr/sbin/xapt} \\
Lateral Movement & SSH service account hijacking \\
Exfiltration & SoftEther VPN Bridge to offshore infrastructure, URL-safe Base64 \\
Polling & 1s initial $\rightarrow$ 5--10min backoff (randomized) \\
Sanitization & Clears 1,000 rows $\times$ A--Z columns on first execution \\
Operational Scope & 53 organizations across 42 countries, active since 2017 \\
ICE Data Exfiltrated & $\sim$9GB internal emails from locally-hosted server \\
\bottomrule
\end{tabularx}

\subsection{Google Sheets Cell Map}

\begin{tabularx}{\textwidth}{lX}
\toprule
\textbf{Cell} & \textbf{Function} \\
\midrule
A1 & Command register — polled for shell commands, S-C-R status on exec \\
A2--An & Data transfer — 45KB fragments for exfil/payload upload \\
V1 & Host fingerprint — hostname, IP, OS, user, language \\
\bottomrule
\end{tabularx}

\subsection{Detection Indicators for ICE SOC}

\begin{tabularx}{\textwidth}{lX}
\toprule
\textbf{Category} & \textbf{Indicator} \\
\midrule
Filesystem & Search for \texttt{/usr/sbin/xapt} binary on all Linux systems \\
Persistence & Audit systemd services for \texttt{xapt.service} \\
Network & Monitor Google Sheets API OAuth tokens in outbound traffic \\
Cryptographic & Check for AES-128-CBC key files (16-byte) in service account directories \\
Exfiltration & Detect SoftEther VPN bridge connections to unknown endpoints \\
Traffic & Scan for URL-safe Base64 encoded payloads in HTTP traffic \\
\bottomrule
\end{tabularx}

\subsection{Known Precedents}

\begin{itemize}
  \item \textbf{TOUGHPROGRESS / TONGUESHED} (2024) — China-nexus; Google Sheets API read/write as tasking/exfil
  \item \textbf{Voldemort Campaign} (Aug 2024, Proofpoint) — Google Sheets C2, AES payloads, DLL side-loading; $\sim$70 organizations across 18 countries
  \item \textbf{T1543.002} — MITRE ATT\&CK technique for systemd service-level persistence
\end{itemize}

\section{Vulnerability Assessment}

\subsection{CVE Chain — 12 Confirmed Vulnerabilities}

\begin{longtable}{p{2.8cm}p{2.5cm}p{5cm}p{2.5cm}}
\toprule
\textbf{CVE} & \textbf{Target} & \textbf{Mechanism} & \textbf{CVSS} \\
\midrule
\endhead
CVE-2022-41328 & FortiOS & Path traversal & — \\
CVE-2023-20867 & VMware Tools & Guest Operations bypass & — \\
CVE-2023-30799 & MikroTik RouterOS & Privilege escalation & — \\
CVE-2023-34048 & VMware vCenter & Out-of-bounds write (DCERPC) & — \\
CVE-2024-5947 & DSE855 Gateway & HTTP GET /Backup.bin $\rightarrow$ plaintext SCADA creds & 6.5 Med \\
CVE-2024-5950 & DSE855 Gateway & Stack buffer overflow $\rightarrow$ RCE on NXP LPC4357 & — \\
CVE-2024-5949 & DSE855 Gateway & Multipart boundary $\rightarrow$ infinite loop DoS & — \\
CVE-2024-5952 & DSE855 Gateway & Unauth web UI restart $\rightarrow$ device reboot & — \\
\rowcolor{lightgray}
CVE-2025-10948 & MikroTik REST API & Heap overflow in libjson.so & \textcolor{critred}{\textbf{8.8--9.8 Crit}} \\
CVE-2025-21590 & Juniper MX & Improper isolation of shared resources & — \\
CVE-2025-9491 & Windows LNK & Zero-day exploitation & — \\
ZDI-CAN-25373 & Windows LNK & Spear-phishing via LNK & — \\
\bottomrule
\end{longtable}

\subsection{SETECOM / SCADA Critical Exposure}

\begin{center}
\colorbox{critred!10}{
\begin{minipage}{0.9\textwidth}
\vspace{0.5em}
\textbf{\textcolor{critred}{CRITICAL: Default Credentials Across Entire Fleet}}\\[0.3em]
\textbf{Username:} \texttt{Admin} \hspace{2cm} \textbf{Password:} \texttt{Password1234}\\[0.3em]
Applied to: DSE 855, DSE 890 MKII, DSE 891, DSE 892 gateways, DSE WebNet cloud management.\\
Controls backup generators for ICE national power grid, Liberty telecom, hospitals, cellular towers.
\vspace{0.5em}
\end{minipage}
}
\end{center}

\begin{tabularx}{\textwidth}{llX}
\toprule
\textbf{Protocol} & \textbf{Port} & \textbf{Vulnerability} \\
\midrule
Modbus TCP/IP & 502 & No authentication, no encryption, no integrity checking by protocol design \\
SNMP v2c & 161/162 UDP & Cleartext community strings ``public'' (read) / ``private'' (read-write) routinely unchanged \\
J1939 CAN Bus & — & 29-bit identifier; direct ECU access for fuel injection, RPM, engine overspeed \\
DSE WebNet & HTTP & Default Admin/Password1234 across entire fleet \\
\bottomrule
\end{tabularx}

\vspace{0.5em}
\textbf{Supply Chain Gap:} The US-funded SOC (Mandiant/Google SecOps) monitors application-layer traffic. The underlying OT infrastructure (DSE controllers, Modbus, SNMP) remains accessible at Tier-0 through SETECOM's architecture. UNC3886 persists at hypervisor level while SOC monitors only above guest OS.

\section{Malware Arsenal}

\begin{tabularx}{\textwidth}{llXXl}
\toprule
\textbf{Family} & \textbf{Layer} & \textbf{Mechanism} & \textbf{Persistence} & \textbf{Actor} \\
\midrule
REPTILE & Kernel & Loadable kernel module & Auto-loads via modprobe.d & UNC3886 \\
MEDUSA & Kernel & Credential harvester via LD\_PRELOAD & Intercepts PAM auth & UNC3886 \\
GRIDTIDE & Application & Google Sheets API C2 & systemd service (xapt) & UNC2814 \\
SOGU.SEC/PlugX & Memory & DLL side-loading, memory-resident & No disk footprint & UNC6384 \\
TINYSHELL & Application & Custom reverse shell variants & Juniper FreeBSD process injection & UNC3886 \\
STATICPLUGIN & Application & Signed downloader via AitM captive portal & Valid code-signing certs & UNC6384 \\
\bottomrule
\end{tabularx}

\vspace{0.5em}
\textbf{Dwell Time:} Median dwell time for UNC3886 espionage operations is \textbf{122 days} ($\sim$10$\times$ average eCrime). Operations occur below guest OS at Tier-0 hypervisor level — invisible to all EDR solutions.

\section{Geopolitical Evidence Chain}

\begin{enumerate}[label=\arabic*.]
  \item \textbf{Budapest Convention 5G Decree (Aug 2023)} — Executive decree mandating all 5G vendors be Budapest Convention signatories, effectively banning Huawei and ZTE from Costa Rican infrastructure
  \item \textbf{ICE Internal Workers' Front Proxy} — Beijing reportedly leveraged ICE's internal union to appeal the 5G ban in lower courts, stalling deployment and creating institutional friction
  \item \textbf{Chinese Embassy Rejection} — Embassy publicly rejected ICE breach accusations, demanded technical evidence from Costa Rican authorities
  \item \textbf{Document 79 Directive} — PRC ordered state-owned enterprises in financial/energy sectors to replace all Western cybersecurity software by 2027
  \item \textbf{Huawei Legal Appeals} — Huawei initiated aggressive legal appeals through Costa Rican Constitutional Court to challenge 5G exclusion
  \item \textbf{US \$25M Cybersecurity Grant} — US State Department committed \$25M to assist Costa Rica in building defenses
  \item \textbf{Supply Chain Gap} — US-funded SOC monitors application layer while underlying edge devices remain PRC-accessible at Tier-0
\end{enumerate}

\section{Prioritized Remediation Plan}

\begin{enumerate}[label=\textbf{\arabic*.}]
  \item \priority{critred}{IMMEDIATE} \textbf{Audit all DSE855/890/891/892 gateways for default credentials}\\
  Every SETECOM-distributed unit ships with Admin/Password1234. These control backup generators for ICE's national grid, hospitals, and cell towers.

  \item \priority{critred}{IMMEDIATE} \textbf{Sweep for GRIDTIDE persistence}\\
  Check all Linux servers for \texttt{/etc/systemd/system/xapt.service} and \texttt{/usr/sbin/xapt}. GRIDTIDE uses Google Sheets API as C2 — traditional network IDS will not detect it.

  \item \priority{highorange}{HIGH} \textbf{Deploy Tier-0 hypervisor monitoring}\\
  UNC3886 operates below guest OS. Current SOC monitors application layer only. REPTILE and MEDUSA rootkits are invisible to EDR.

  \item \priority{highorange}{HIGH} \textbf{Patch MikroTik fleet against CVE-2025-10948}\\
  Critical heap overflow (CVSS 8.8--9.8) in REST API. MikroTik routers are ubiquitous in Costa Rican ISP infrastructure.

  \item \priority{highorange}{HIGH} \textbf{Investigate TR-069 management plane access}\\
  Unauthorized admin password reset detected 2026-01-30 via TR-069 on ARRIS router. This protocol gives ISPs full remote control over CPE.

  \item \priority{medyellow}{MEDIUM} \textbf{Audit Modbus TCP/502 exposure}\\
  Modbus has zero authentication by design. Segment and firewall port 502 on all SCADA networks.

  \item \priority{medyellow}{MEDIUM} \textbf{Review SoftEther VPN connections}\\
  GRIDTIDE uses SoftEther VPN Bridge for exfiltration. Check for unexpected outbound VPN tunnels.

  \item \priority{stratblue}{STRATEGIC} \textbf{Engage KAPPA platform for continuous monitoring}\\
  24/7 autonomous correlation across RF, satellite, network, and acoustic domains with automated evidence chain.
\end{enumerate}

\section{KAPPA Platform Capabilities}

The KAPPA platform currently provides:

\begin{itemize}
  \item RF spectrum monitoring via 33 KiwiSDR nodes (71 frequency targets)
  \item Satellite constellation tracking (Beidou, Yaogan, Gaofen ISR)
  \item Network forensics — TR-069, PCAP analysis, MAC fingerprinting
  \item External data feeds — USGS seismic, NOAA space weather, WWLLN lightning
  \item Automated evidence chain with SHA-256 integrity hashing
  \item Device fleet monitoring with latency, jitter, and health scoring
  \item Real-time threat scoring ($\kappa$ score, 0--100 scale)
  \item Interactive intelligence briefing at \url{https://ciajw.com}
\end{itemize}

\section{Contact}

\begin{center}
\begin{tabular}{ll}
\textbf{Name:} & Samuel Wotton \\
\textbf{Project:} & KAPPA SIGINT / Project Echo \\
\textbf{Web:} & ciajw.com \\
\textbf{Location:} & Aparthotel Suites Cristina, Sabana Norte, San Jos\'e, Costa Rica \\
\textbf{Coordinates:} & 10.0514\textdegree N, 84.2187\textdegree W \\
\end{tabular}
\end{center}

\vspace{1em}
\begin{center}
\textit{I'm right next door. If your team would like to discuss any of these findings,\\review the raw data, or explore how KAPPA can help — I'm available anytime.}
\end{center}

\vfill
\begin{center}
\rule{0.5\textwidth}{0.4pt}\\[0.5em]
{\small KAPPA SIGINT Platform v1.0 — \today}\\
{\small All findings derived from passive collection and open-source intelligence.}\\
{\small No ICE systems were accessed or tested.}
\end{center}

\end{document}
